
\documentclass[a4paper,12pt]{article}
\usepackage{amsmath,graphicx,hyperref}
\usepackage{xcolor}

\title{\textbf{An Exploration of Random LaTeX Structures}}
\author{Hinata Research Group}
\date{\today}

\begin{document}
\maketitle

\begin{abstract}
This document demonstrates a random assortment of \LaTeX{} features such as equations, tables, and figures. It serves as a playground for document generation and formatting tests.
\end{abstract}

\section{Introduction}
LaTeX is widely used for academic and scientific writing because it handles mathematical typesetting, references, and document layout elegantly.

\section{Mathematical Expression}
Consider the quadratic formula:
\[
x = \frac{-b \pm \sqrt{b^2 - 4ac}}{2a}
\]
It describes the roots of any quadratic equation \( ax^2 + bx + c = 0 \).

\subsection{Matrix Example}
An example of a 3x3 matrix:
\[
A = 
\begin{bmatrix}
1 & 2 & 3 \\
4 & 5 & 6 \\
7 & 8 & 9
\end{bmatrix}
\]

\section{Table Example}
\begin{center}
\begin{tabular}{|c|c|c|}
\hline
\textbf{Variable} & \textbf{Value} & \textbf{Description} \\
\hline
$a$ & 2 & Coefficient of $x^2$ \\
$b$ & 5 & Coefficient of $x$ \\
$c$ & 3 & Constant term \\
\hline
\end{tabular}
\end{center}

\section{Figure Example}
Figure~\ref{fig:example} shows a placeholder image.

\begin{figure}[h!]
\centering
\includegraphics[width=0.5\textwidth]{example-image}
\caption{An example placeholder image from the \texttt{graphicx} package.}
\label{fig:example}
\end{figure}

\section{Conclusion}
This was a random LaTeX document demonstrating multiple environments. For more, see the \href{https://www.overleaf.com/learn}{Overleaf documentation}.

\end{document}
